\documentclass[11pt]{article}

\usepackage[a4paper, margin=1in]{geometry}
\usepackage[T1]{fontenc}
\usepackage[utf8]{inputenc}
\usepackage{lmodern}
\usepackage{microtype}
\usepackage[parfill]{parskip}
\usepackage{xcolor}
\usepackage{titlesec}
\usepackage{titling}
\usepackage{enumitem}
\usepackage[hidelinks]{hyperref}
\usepackage{fancyhdr}

\definecolor{redesteal}{HTML}{047F78}
\definecolor{redesdark}{HTML}{242B2E}
\definecolor{redesneutral}{HTML}{6E6E6E}

\titleformat{\section}
  {\normalfont\Large\bfseries\color{redesdark}}
  {}{0pt}{}
\titlespacing*{\section}{0pt}{1.4em}{0.4em}

\pagestyle{fancy}
\fancyhf{}
\fancyfoot[L]{\small\color{redesneutral}REDES-IA \textemdash{} Formación}
\fancyfoot[R]{\small\color{redesneutral}\thepage}
\renewcommand{\headrulewidth}{0pt}
\renewcommand{\footrulewidth}{0pt}

\begin{document}

\thispagestyle{fancy}

\noindent{\color{redesteal}\rule{\linewidth}{1.5pt}}\par
\vspace{0.5em}
\noindent{\Huge\bfseries\color{redesdark}\hyphenpenalty=10000\exhyphenpenalty=10000 Agentic Coding for Academic Researchers\par}
\vspace{0.5em}
\noindent{\large\color{redesneutral}Universitat de Barcelona \textemdash{} Departament d'Economia\par}
\vspace{0.4em}
\noindent{\color{redesteal}\rule{\linewidth}{0.5pt}}\par
\vspace{1.4em}

\noindent{\bfseries Instructor:} Brady Allardice \\
\noindent{\bfseries Sessions:} 1, 3, 5 and 8 June 2026 (four sessions, eight hours) \\
\noindent{\bfseries Schedule:} Mon 1 June, 15.00\,--\,17.00 \textbullet{} Wed 3 June, 15.00\,--\,17.00 \textbullet{} Fri 5 June, 10.00\,--\,12.00 \textbullet{} Mon 8 June, 15.00\,--\,17.00

\section*{Course description}

A hands-on workshop (approximately 60\,--\,70\% exercises) that teaches quantitative researchers how to use Claude Code, an agentic AI coding tool, to accelerate their research workflows. Participants learn how to harness agentic coding to clean data, generate robustness checks, audit code, and build replication packages, while maintaining full analytical control over every decision. The course covers effective instruction-writing, execution modes, version control with Git, failure diagnosis, and the Anthropic API for batch processing. No prior experience with AI coding tools is required; comfort with quantitative research and data analysis in any language is sufficient.

\section*{Audience and prerequisites}

For grad students, postdocs, and faculty across economics and the social sciences. Comfort with one statistical environment (Python, R, or Stata) is assumed; no prior agent or LLM-tooling experience required. Before Session 1, participants install Claude Code CLI, Python 3.9+ with pandas / numpy / statsmodels / matplotlib, git, VS Code with the Claude Code and LaTeX Workshop extensions, and a TeX distribution. A prework setup guide is distributed one week in advance.

\section*{Learning outcomes}

By the end of the workshop, participants will be able to: distinguish agentic AI from chatbots and autocomplete tools, and pick the right approach for the task at hand; manage context windows, scoping, and verification across long agentic sessions; use git, GitHub, and LaTeX in VS Code as one integrated environment alongside Claude Code; build custom skills, subagents, hooks, and MCP connections to extend the agent for recurring research tasks; and replicate a published empirical paper from raw data to final figure with agentic AI as the primary tool, while retaining methodological control.

\section*{Schedule}

\textbf{Session 1 (Mon 1 June) — Introduction to Claude Code.} What Claude Code is and where it sits in the spectrum from manual coding through autocomplete, web chat, and full agentic systems. The RA analogy: delegating concrete tasks, verifying outputs, retaining control. Terminal, VS Code, scripts vs.\ notebooks, virtual environments, context windows, and permission modes. Exercise: set up the project, run a small data task end to end, and write a first CLAUDE.md.

\textbf{Session 2 (Wed 3 June) — Research pipelines.} Working within the context window: five practical strategies for keeping long sessions coherent. The core agentic workflow — describe, review, generate and validate, save. Six validation strategies for catching agent errors before they propagate. Exercise: the Swiss Franc Shock data pack — merge, explore, specify, estimate, run robustness checks, and produce a publication-ready LaTeX regression table.

\textbf{Session 3 (Fri 5 June) — Git and LaTeX in VS Code.} Why version control matters more, not less, when an agent edits your files directly. Minimum-viable git: init, add, commit, diff, log, checkout. \texttt{.gitignore} and \texttt{.claudeignore}. Moving from Overleaf to LaTeX in VS Code so paper, code, and data all live in one project. Exercise: build a short paper from scratch using git, LaTeX, and Claude Code together.

\textbf{Session 4 (Mon 8 June) — Skills, MCPs, subagents, and capstone.} Skills as packaged knowledge for recurring tasks (abstract audits, regression tables, spec validation) and the routing rule for where they belong. MCPs as connections to external tools, with the Zotero MCP for literature work and the GitHub MCP for collaboration. Subagents as roles (a referee, a copyeditor, a methodologist) versus skills as knowledge; hooks for automation (research audit logs, auto-commit). Capstone: replicate Figure~1 of a published empirical paper from raw data, run a small extension, and write up the result.

\section*{Assessment}

No grades. Two tracks for the capstone: (a) a guided replication of a provided paper, or (b) integration of Claude Code into one stage of your own ongoing project, using at minimum a CLAUDE.md, plan mode, git, and one power feature (hook, skill, MCP, or subagent). Deliverable is a version-controlled project folder containing data, code, output, and a compiled .tex write-up.

\section*{Readings}

Allardice, B. (2026). \emph{Ten Ways to Use Agentic AI in Academic Research}. Course notes (distributed in advance). \\
Anthropic. \emph{Claude Code documentation}. \href{https://docs.claude.com/en/docs/claude-code}{docs.claude.com/en/docs/claude-code} \\
Anthropic. \emph{Building Effective AI Agents}. \href{https://www.anthropic.com/engineering/building-effective-agents}{anthropic.com/engineering/building-effective-agents} \\
Course repository and exercise materials (distributed at enrollment).

\section*{Instructor}

Brady Allardice is a PhD candidate in Political and Social Science at Universitat Pompeu Fabra (UPF) in Barcelona. His research sits at the intersection of artificial intelligence, technological change, and politics. His thesis examines how automation and AI shape support for the political establishment — and how political institutions, in turn, steer the direction of technological change. He is supervised by Aina Gallego. His background is in Economics and quantitative methods, having worked as an Economic Consultant and a Data Scientist after completing a Master in Economics at the Barcelona School of Economics.

\end{document}
